\subsection{A genus-four adjacent descending link in the source argument}

\begin{result}{D.72}{Refuted}{The genus-four adjacent descending-link claim}
For the cycle complex and coefficient-sum weight function of
\cite{D_Minahan2025}, there is a
weight-two constant edge $\sigma$ on a closed genus-four surface
whose actual adjacent descending link is an infinite discrete set.
In particular,
\[
 \widetilde H_0\bigl(d_W^{\rm adj}(\sigma);\mathbb Z\bigr)\ne0,
 \qquad g-3-\dim\sigma=0.
 \tag{D.600}
\]
This refutes the adjacent-descending-link bound used in the
final proof on page~29.  The refuted assertion is this local
bound for the specified function and complex.
The construction does not produce a nonzero homology class in
the full complex of homologous curves and does not decide
its claimed acyclicity by another argument.
The facial descending link and the maps to the printed
global convex join are also computed below, so no
equivalence between those variants is assumed.
\end{result}

\begin{proof}
\medskip\noindent\textbf{The gluing maps and the homology detector.}
Let $P=S_{2,2}$ and $Q=S_{0,4}$ be connected compact oriented
surfaces.  Give their boundary components labels
\[
 \partial P=P_a\sqcup P_b,\qquad
 \partial Q=Q_a\sqcup Q_b\sqcup Q_{c,+}\sqcup Q_{c,-}.
 \tag{D.601}
\]
Choose orientation-compatible circle parametrizations of
these boundary components.  Make the actual identifications
\[
 \begin{aligned}
 P_a(e^{it})&\sim Q_a(e^{-it}),&
 P_b(e^{it})&\sim Q_b(e^{-it}),\\
 Q_{c,+}(e^{it})&\sim Q_{c,-}(e^{-it}).
 \end{aligned}
 \tag{D.602}
\]
The resulting closed oriented surface $S$ has
$\chi(S)=\chi(P)+\chi(Q)=-4-2=-6$, hence genus four.
Denote the three seams by $a,b,c$.  They are essential
and pairwise nonisotopic.  The curves $a,b$ bound $P$;
the other side is the genus-one surface with two boundaries
obtained by gluing $Q_{c,+}$ to $Q_{c,-}$.
Thus $a,b$ are not parallel.  Cutting $a$ and $c$ leaves
a connected surface, since the retained $b$-gluing joins
$P$ to $Q$.  Consequently $a,c$ are a partial cut system.

Orient $a,b,c$ so that
\[
 [a]=[b]=v,\qquad [c]=w,\qquad x=v+w,
 \tag{D.603}
\]
where $v,w$ are independent and extend to a primitive
isotropic basis of $H_1(S;\mathbb Z)$.
With these choices, the induced boundary classes of $Q$ are
\[
 \begin{array}{c|rrrr}
 \text{boundary}&Q_a&Q_b&Q_{c,+}&Q_{c,-}\\ \hline
 \text{class in }H_1(S;\mathbb Z)&-v&v&w&-w .
 \end{array}
 \tag{D.604}
\]
Thus $Q_b,Q_{c,+}$ are the incoming pair and
$Q_a,Q_{c,-}$ the outgoing pair, after fixing the
source's cobordism convention.  The complementary
piece $P$ has one incoming and one outgoing boundary.
Both pieces are balanced.

There is an arc in $Q$ from $Q_{c,-}$ to $Q_{c,+}$,
away from its other boundaries.  Close it across the
$c$-gluing to obtain a closed curve $\gamma$.
It is disjoint from $P,a,b$ and meets $c$ once.
Its orientation can be chosen so that the actual
intersection homomorphism satisfies
\[
 \langle w,[\gamma]\rangle=1,\qquad
 \langle v,[\gamma]\rangle=0,\qquad
 \langle x,[\gamma]\rangle=1.
 \tag{D.605}
\]
In particular $x$ is primitive, and no curve contained
in $P$ represents $x$.

\medskip\noindent\textbf{An actual cell, with its full coefficient polytope.}
In the source, a basic cycle for $x$ is an oriented
multicurve admitting a unique positive integral
coefficient expression for $x$.
A support $M$ is a cycle when each of its curves occurs
in such a basic cycle contained in $M$ and no nonzero
nonnegative integral combination of its curve classes
vanishes.  Its cell $P_M$ is the convex hull of its
basic cycles in the vector space with curve classes
as formal coordinate labels.

For $M=a\sqcup b\sqcup c$, the complete coefficient
polytope is
\[
 P_M=\{(\lambda_a,\lambda_b,\lambda_c)\in
       \mathbb R_{\ge0}^3:
       \lambda_a+\lambda_b=1,\ \lambda_c=1\}.
 \tag{D.606}
\]
Its only vertices are $(1,0,1)$ and $(0,1,1)$.
They are the basic cycles $a+c$ and $b+c$.
Every support curve occurs in one of them.
Independence of $v,w$ implies
\[
 r[a]+s[b]+t[c]=0,\quad r,s,t\ge0
       \quad\Longrightarrow\quad r=s=t=0.
 \tag{D.607}
\]
Thus both defining cycle conditions hold.
The segment $\sigma=P_M$ is an actual one-cell of the
specified complex.  The source's function $W$ sums
the curve coefficients, so it is identically two on
$\sigma$.

\medskip\noindent\textbf{Classification of every lower vertex.}
Retain the source's literal set $R_\sigma$ of cells
$\rho$ such that $\sigma$ is a face of $\rho$ and
\[
 W(u)\le2\quad(u\in\rho^{(0)}),\qquad
 W(u)=2\quad\Longleftrightarrow\quad u\in\sigma^{(0)}.
 \tag{D.608}
\]
Every vertex outside $\sigma$ has positive integral
weight less than two.  Hence it is a single oriented
curve $\delta$ representing $x$, with coefficient one.
The support of $\rho$ is a multicurve containing
$a,b,c$, so $\delta$ is disjoint from the seams.
It cannot be a seam and must lie in the interior of
$P$ or $Q$.  Equation~(D.605) excludes $P$.

An essential nonperipheral simple closed curve in $Q$
separates its four boundary components into two pairs.
Summing (D.604) over one side gives the three possible
unordered partition classes, up to sign:
\[
 \begin{array}{c|c}
 \text{partition}&\text{curve class up to orientation}\\ \hline
 \{a,b\}\mid\{c_+,c_-\}&0\\
 \{a,c_+\}\mid\{b,c_-\}&v-w\\
 \{a,c_-\}\mid\{b,c_+\}&v+w .
 \end{array}
 \tag{D.609}
\]
Peripheral curves do not represent $x$ either.
Thus every lower vertex is a curve with the incoming
pair on one side and the outgoing pair on the other.
Conversely, orienting any such curve appropriately
gives class $x$.  Write $\mathscr D$ for this set of
oriented isotopy classes.

\medskip\noindent\textbf{Every classified vertex has a descending triangle.}
For $\delta\in\mathscr D$, the support $M\sqcup\delta$
has coefficient polytope
\[
 \begin{gathered}
 \lambda_a+\lambda_b+\lambda_\delta=1,\qquad
 \lambda_c+\lambda_\delta=1,\qquad
 \lambda_a,\lambda_b,\lambda_c,\lambda_\delta\ge0,\\
 P_{M\sqcup\delta}
    =\operatorname{conv}\{(1,0,1,0),(0,1,1,0),(0,0,0,1)\}.
 \end{gathered}
 \tag{D.610}
\]
These three vertices are precisely $a+c,b+c,\delta$,
and each is basic.  Every support curve occurs.
The linear functional taking both $v$ and $w$ to one
is positive on every support class, including
$[\delta]=v+w$.  It therefore excludes every
nonzero nonnegative null combination.  This verifies
the source's cycle conditions for the entire support.
The triangle is consequently an actual cell.

On this triangle the weight is
\[
              W=2-\lambda_\delta.
 \tag{D.611}
\]
Its exact maximum face is $\sigma$, and its remaining
vertex has weight one.  It belongs to $R_\sigma$;
$\delta$ is adjacent to both vertices of $\sigma$.
This proves that every element of $\mathscr D$ occurs
in the actual adjacent descending link.

\medskip\noindent\textbf{No cell connects distinct lower vertices.}
Two distinct essential nonperipheral curves in a
four-holed sphere cannot be disjoint: cutting along
one leaves two pairs of pants, whose essential closed
curves are boundary-parallel.  These classes and their
intersection numbers remain distinct after (D.602).
Indeed a bigon in the closed surface with boundary
on the two curves can be reduced to one in $Q$.
Its boundary is disjoint from the seams, and no
essential seam can be contained in its disk; the
usual innermost-circle removal puts the disk in $Q$.
The bigon criterion therefore computes the same
minimal intersection in $Q$ and in $S$.

For a concrete infinite family, fix
$\delta_0\in\mathscr D$ and choose $\eta\subset Q$
with partition $\{a,b\}\mid\{c_+,c_-\}$ and
$i(\eta,\delta_0)=2$.  Its class in $S$ is zero.
In fact its two complementary components in $S$
have genera three and one.  The actual twist
$T_\eta$ fixes every boundary component of $Q$ and
preserves the partition of $\delta_0$.
Set $\delta_m=T_\eta^m(\delta_0)$.  Twisting the
two crossing arcs in an annular collar of $\eta$
gives
\[
  \delta_m\in\mathscr D,\qquad
  i(\delta_m,\delta_n)
       =|m-n|\,i(\eta,\delta_0)^2=4|m-n|.
 \tag{D.612}
\]
For this intersection calculation, the two arcs of
each curve cross the collar from one boundary to
the other.  Lifting the collar to a strip makes
the relative winding $m-n$ explicit; each pair
of crossing arcs contributes $|m-n|$ intersections.
Outside the collar the curves are parallel copies.
The bigon criterion leaves exactly these
$4|m-n|$ intersections.  In particular
$\delta_0,\delta_1$ are distinct actual lower vertices.

A support multicurve of a cell $\rho\in R_\sigma$
can contain at most one member of $\mathscr D$.
Every vertex of $\rho$ has already been classified
by (D.608).  The cycle support is the union of its
basic-cycle supports, so $\rho$ is either $\sigma$
itself or one of the triangles (D.610).
If $V(\rho)$ denotes the source's set of vertices
outside $\sigma$ adjacent to a vertex of $\sigma$,
each nonempty $\operatorname{Hull}V(\rho)$ is
therefore a singleton.  We have the exact
identifications, with their actual vertex maps,
\[
 d_W^{\rm adj}(\sigma)
       =\coprod_{\delta\in\mathscr D}\{\delta\},
 \qquad
 d_W^{\rm face}(\sigma)
       =\coprod_{\delta\in\mathscr D}\{\delta\}.
 \tag{D.613}
\]
For the facial equality, the only face of a triangle
(D.610) disjoint from $\sigma$ is its vertex $\delta$.
The difference $[\delta_1]-[\delta_0]$ is thus a
nonzero reduced zero-homology class in both links.
This proves (D.600).

\medskip\noindent\textbf{The exact source application and its variants.}
Lemma~6.1 on printed page~22 assumes
$|B_+|\ge2$ and $|B_-|\ge2$.
The balanced component $P=S_{2,2}$ has one of each.
Its draining complex is empty: a draining subsurface
would have at most one incoming boundary and at
least one new outgoing boundary, contradicting the
strict draining inequality.  The factor for $Q$ is
the discrete set (D.613).  The main proof's invocation
of that lemma for every complementary component
therefore does not provide its asserted
$(g-3-\dim\sigma)=0$ acyclicity in this example.
The calculation above verifies the failure directly,
without relying on the source's Lemma~6.6.

The printed total-link definition on page~15 is
the \emph{global} convex join
\[
 d_W^{\rm tot}(\sigma)=
   \operatorname{CJoin}
       \bigl(d_W^{\rm face}(\sigma),d_W^{\rm adj}(\sigma)\bigr).
 \tag{D.614}
\]
With its displayed definition of convex join, (D.613)
makes this the union of all line segments between
the formal coordinate vectors labelled by
$\mathscr D$.  In particular it contains
$[\delta_0,\delta_1]$, which is unsupported by any
cell of the cycle complex: its endpoints intersect.
The actual inclusion from either link in (D.613)
to (D.614) sends
$[\delta_1]-[\delta_0]$ to zero in reduced $H_0$.
It cannot be a weak homotopy equivalence.
Thus the printed global-join formula and the
equivalences asserted in Lemma~5.3 do not supply
a repair of the adjacent-link bound.  This calculation
also explains why the conclusion of this result is
stated for the adjacent link, without silently
substituting a different cellwise definition of the
total link.

\medskip\noindent\textbf{Source and certificate scope.}
The official arXiv record was reopened on
16~November~2026 and still listed version~3,
4~April~2025, as its latest version.
The author's current homepage and CV were checked;
the linked May~2024 thesis has the same draining
hypotheses and final application.
No later repair was located by these checks.
The retrieved definitions and the complete final
application were read; full-size renders of printed
pages~22 and~29 were also inspected.

The exact finite gluing, homology-partition and
polytope checks are retained in
\path{code/d_minahan_constant_edge.py}, its JSON,
stdout and stderr.  All nineteen checks passed.
The first execution exposed a one-use-generator
coding error in the vector-sum helper.  Its original
code, empty stdout and traceback are preserved;
materializing the input once repaired that helper,
and the complete rerun produced the stated certificate.
The curve classification and intersection arguments
are the proofs above; the program does not purport
to enumerate their infinite isotopy classes.

The proposed ambient-Torelli finite-homology derivation
is retained as an unpromoted research note, with
this precise source-input obstruction.
Its independent arithmetic and module calculations
are unchanged.  No finiteness theorem is substituted
for the missing homological comparison, and no
promoted result elsewhere in Part~D depends on
this source input.
\end{proof}
